% !TEX root = ../Main.tex
\chapter{冲上冲回的速度比；斜面加斜向力}

接上一章：物块冲上斜面、再滑回原处，回来时的速度总比冲上去时小，差的那一截被摩擦力耗掉了。把上行、下行两段用动能定理写出来一除，就得到一个干净的速度比。最后看一类容易错的题——斜面上加一个斜向的外力，此时 $N\ne mg\cos\theta$。

\section{冲上、冲回的速度比}

物块以 $v_1$ 沿斜面冲上，到顶后滑回原处时速度为 $v_2$（要能滑回，须 $\mu<\tan\theta$）。设上行、下行的距离都是 $L$。

\begin{center}
\begin{tikzpicture}[>=Stealth, scale=1]
    \draw[thick] (0,0) -- (5,2.5);
    \draw[thick] (0,0) -- (5,0);
    \draw (0.7,0) arc (0:27:0.7);
    \node at (0.95,0.18) {\small $\theta$};
    \node at (4.3,2.0) {\small $\mu$};
    \coordinate (C) at (2.0,1.0);
    \draw[thick, fill=gray!15] (C) ++(207:0.3) -- ++(27:0.6) -- ++(117:0.46) -- ++(27:-0.6) -- cycle;
    \draw[->, very thick, blue] (C) ++(27:0.35) -- ++(27:1.0) node[right] {\small $v_1$};
    \draw[->, very thick, red] (C) ++(207:0.35) -- ++(207:1.0) node[below left] {\small $v_2$};
\end{tikzpicture}
\end{center}

上行用动摩擦力向下、下行用动摩擦力向上，对两段分别写动能定理：
\[
\begin{aligned}
\text{冲上：}\quad & 0-\tfrac12 m v_1^2=-mg\sin\theta\cdot L-\mu mg\cos\theta\cdot L,\\
\text{冲回：}\quad & \tfrac12 m v_2^2-0=mg\sin\theta\cdot L-\mu mg\cos\theta\cdot L.
\end{aligned}
\]
两式相除（$L$、$\tfrac12 m$ 都约掉）：
\[
\pr{\frac{v_1}{v_2}}^2=\frac{\sin\theta+\mu\cos\theta}{\sin\theta-\mu\cos\theta}.
\]
分母小于分子，所以 $v_1>v_2$——滑回来一定比冲上去慢。

\section{相似问题（引申）：多次往返}

如果物块在两个斜面（或斜面与某种反弹）之间一次次往返，每一个"上去再下来"的循环几何完全一样。

\begin{center}
\begin{tikzpicture}[>=Stealth, scale=1]
    % V 形谷：左右两个斜面在底部相交
    \draw[thick] (-2.6,1.7) -- (0,0) -- (2.6,1.7);
    \draw[thick] (-2.6,0) -- (2.6,0);
    % 右斜面上的物块，速度 v1 向右上
    \coordinate (R) at (1.15,0.75);
    \draw[thick, fill=gray!15] (R) ++(213:0.28) -- ++(33:0.56) -- ++(123:0.44) -- ++(33:-0.56) -- cycle;
    \draw[->, very thick, blue] (R) ++(33:0.3) -- ++(33:1.0) node[right] {\small $v_1$};
    % 左斜面上的物块，速度 v2 向左上
    \coordinate (L) at (-1.15,0.75);
    \draw[thick, fill=gray!15] (L) ++(327:0.28) -- ++(147:0.56) -- ++(57:0.44) -- ++(147:-0.56) -- cycle;
    \draw[->, very thick, red] (L) ++(147:0.3) -- ++(147:1.0) node[left] {\small $v_2$};
\end{tikzpicture}
\end{center}

相邻两次的速度比是同一个常数：
\[
\pr{\frac{v_{n-1}}{v_n}}^2=\frac{\sin\theta+\mu\cos\theta}{\sin\theta-\mu\cos\theta}.
\]
也就是说，每往返一次速度按固定比例衰减，是个等比数列。

\section{斜面加斜向力}

物块在倾角 $\theta$ 的斜面上，受一个\textbf{水平方向}的外力 $F$（把它压向斜面）。取 $x$ 轴沿斜面向上、$y$ 轴沿斜面法向，把 $F$ 和 $mg$ 都分解到这两个方向。

\begin{center}
\begin{tikzpicture}[>=Stealth, scale=1]
    \draw[thick] (0,0) -- (5,2.5);
    \draw[thick] (0,0) -- (5,0);
    \draw (0.7,0) arc (0:27:0.7);
    \node at (0.95,0.18) {\small $\theta$};
    \coordinate (C) at (2.4,1.2);
    \draw[thick, fill=gray!15] (C) ++(207:0.32) -- ++(27:0.64) -- ++(117:0.5) -- ++(27:-0.64) -- cycle;
    \draw[->, very thick, blue] (C) ++(180:1.3) -- (C) node[midway, above] {\small $F$};
    \draw[->, thick] (C) -- ++(27:1.3) node[right] {\small $x$};
    \draw[->, thick] (C) -- ++(117:1.3) node[above] {\small $y$};
\end{tikzpicture}
\end{center}

\state{受力分析}{
    法向（$y$）平衡：水平力 $F$ 把物块压向斜面，法向分量是 $F\sin\theta$，与重力法向分量同侧，所以
    \[
    y:\quad N=mg\cos\theta+F\sin\theta.
    \]
    沿斜面（$x$）由牛顿第二定律：
    \[
    x:\quad ma=F\cos\theta-\mu N.
    \]
}

\rmkb[注意]{
    此时 $N\ne mg\cos\theta$！\,只要斜面上有不沿斜面的外力（水平力、竖直力等），它的法向分量就会改变支持力 $N$，进而改变摩擦力 $\mu N$。一定先用 $y$ 方向解出 $N$，再代回 $x$ 方向，别一上来就写 $N=mg\cos\theta$。
}
